\documentclass[conference]{IEEEtran}
\IEEEoverridecommandlockouts
\usepackage{cite}
\usepackage{amsmath,amssymb,amsfonts}
\usepackage{graphicx}
\graphicspath{{../figs/}}
\usepackage{textcomp}
\usepackage{xcolor}
\usepackage{booktabs}
\usepackage{url}
\usepackage{amsthm}
\newtheorem{proposition}{Proposition}
\renewcommand{\topfraction}{0.95}
\renewcommand{\bottomfraction}{0.95}
\renewcommand{\textfraction}{0.05}
\renewcommand{\floatpagefraction}{0.9}
\renewcommand{\dbltopfraction}{0.95}
\renewcommand{\dblfloatpagefraction}{0.9}
\def\BibTeX{{\rm B\kern-.05em{\sc i\kern-.025em b}\kern-.08em
    T\kern-.166em\lower.7ex\hbox{E}\kern-.125emX}}

\begin{document}

\title{Clause-Activation Signatures (CAS) for Attack-Family Attribution in IoMT}

\author{Rakesh Reddy Yakkati, Per-Arne Andersen, Lei Jiao, Ole-Christoffer Granmo, and Rishad Shafik
\thanks{This work was supported by the \textit{SecureIoTM---Ultra-low-energy IoT Intrusion Detection Systems using Logic-based Tsetlin Machines} project (Grant Number 342167), funded by the Research Council of Norway under the TEKNOLOGIKONVER programme. The authors acknowledge the support from the Centre of Artificial Intelligence Research at the University of Agder, Sørlandet Hospital Trust, and Newcastle University for their collaboration in this interdisciplinary research initiative.}% <-this % stops a space
\thanks{Rakesh Reddy Yakkati, Lei Jiao, and Ole-Christoffer Granmo are with the Department of ICT, University of Agder, Grimstad 4879, Norway (emails rakeshry@uia.no, lei.jiao@uia.no, and ole.granmo@uia.no).}
\thanks{Per-Arne Andersen is with the Department of Information Systems, University of Agder, Kristiansand 4630 (email per.andersen@uia.no).}
\thanks{Rishad Shafik is with the Department of Microelectronic Systems, Newcastle University, Newcastle upon Tyne, UK (email rishad.shafik@newcastle.ac.uk).}}

\maketitle

\begin{abstract}
When an intrusion detector flags traffic in the Internet of Medical Things
(IoMT), a label alone is not enough --- the analyst must also know
\emph{which} attack family is at play and \emph{why}. This paper introduces
Clause-Activation Signatures (CAS), a method that turns the clauses of a
trained weighted Tsetlin Machine (TM) into a compact, auditable signature
for each attack family, so that attribution becomes signature matching
rather than an opaque vote. Each signature is extracted with
complementary-pairs stability selection (CPSS) over clause activations,
split by coefficient sign into inculpatory and exculpatory evidence, and
certified by a Shah--Samworth bound on spurious inclusions. Evaluated on
CICIoMT2024 (WiFi/MQTT, family level) across three seeded TMs under a
leak-free fit/selection/report split of the test set, CAS improves over
the raw weighted-vote TM it is extracted from: macro-F1
rises from $0.692 \pm 0.004$ to $0.765 \pm 0.029$, using only 28--53
sign-positive clauses per family from a 2{,}400-clause pool. Sorting
clauses by weight and discarding the weakest $84\%$ beforehand costs the
base detector nothing, lifts CAS to $0.787 \pm 0.023$, and removes its
sensitivity to the stability threshold. The six signatures provably form a
maximally 5-disjunct code on every seed, so any mixture of up to five
families decodes exactly; under 5\% bit-flip noise a thresholded decoder
still recovers 99.8\%--100\%. A binary normal-versus-attack reduction
lifts AUROC from 0.817 to 0.938. Two negative results are reported: weak
open-set novelty detection under correct conformal false-flag control, and
signature identity that does not survive retraining. CAS is a certified
attribution layer that denoises a mediocre base detector, not a
replacement for one.
\end{abstract}

\begin{IEEEkeywords}
attack attribution, disjunct codes, intrusion detection, interpretable
machine learning, IoMT, open-set recognition, stability selection,
Tsetlin Machine
\end{IEEEkeywords}

\section{Introduction and Related Work}
An IoMT intrusion detector must both raise a correct alarm and name the
attack family behind it in a form an analyst can audit, since containment
and clinical risk differ by family. Post-hoc explainers such as SHAP, LIME
and Anchors~\cite{xaiids2025} serve only the first: they explain one
prediction at a time with no stability guarantee. A per-family
\emph{signature} --- interpretable by design, stability-vetted,
attached to a family rather than a flow --- is absent from the
literature.

The Tsetlin Machine (TM)~\cite{granmo2018} is a natural substrate: its
unit of computation is a conjunctive clause over Boolean literals that
stays human-readable after training, and the weighted
TM~\cite{phoulady2020} adds an integer vote weight per clause. TM-based
intrusion detection is established --- rules on
NSL-KDD~\cite{abeyrathna2020}, TMs on CICIoMT2024 and
MedSec-25~\cite{jaiswal2026iomt,jaiswal2026medsec} --- but these explain single
predictions through whichever clauses fired, without error control. Clause-Activation Signatures (CAS) close that gap: per family,
complementary-pairs stability selection
(CPSS)~\cite{meinshausen2010,shah2013} over clause activations keeps the
stably predictive clauses, splits them by coefficient sign into
inculpatory and exculpatory evidence, and carries a Shah--Samworth
certificate on spurious inclusions. Identifiability under simultaneous families follows from
$t$-disjunct matrices~\cite{du1993}. We evaluate on
CICIoMT2024~\cite{dadkhah2024}, the only peer-reviewed multi-protocol IoMT
benchmark meeting standard validity
criteria~\cite{sommer2010,engelen2021,arp2022}; its schema carries no IP or
port fields.

Contributions: a seed-reproducible CAS implementation on TMU~\cite{tmu}
clause activations with a load-bearing sign filter
(\S\ref{sec:cpss}, \S\ref{sec:validation}); three-seed
$\pi_{\mathrm{thr}}$/$B$/$\kappa$ sensitivity maps (\S\ref{sec:sweep});
certified per-family evidence
lists (\S\ref{sec:signatures}); a clause-pruning program that discards
$84\%$ of the pool and improves CAS (\S\ref{sec:prune},
\S\ref{sec:pruned}); a maximal disjunct-code certificate
(\S\ref{sec:disjunct}); a binary reduction (\S\ref{sec:binary}); and two
negative results (\S\ref{sec:taskb}, \S\ref{sec:stability}).

\section{Method: Clause-Activation Signatures}
\subsection{Overview}
Fig.~\ref{fig:cpss} shows the pipeline. Raw flows are first binarised with
the standard thermometer binariser of the TMU library, and a weighted TM
is trained on the resulting Boolean vectors. Everything afterwards works
on the trained model's learned rules: passing flows through the TM yields
clause activations, and CPSS extracts a stable signature per class from
them. Two guarantees accompany the signatures: a statistical
error-control bound on spurious inclusions (\S\ref{sec:cert}) and a
combinatorial identifiability guarantee for mixtures
(\S\ref{sec:disjunct}).

\begin{figure*}[!t]
\centering
\includegraphics[width=\textwidth]{fig_cpss_protocol.pdf}
\caption{CAS extraction and evaluation under the three-way protocol.
The training partition determines the binarizer and weighted TM.
Optional weight pruning precedes CPSS on the fit half; the selection
quarter determines $k^*$, $B$, and $\pi_{\mathrm{thr}}$ (dashed arrows).
Each complementary half contributes one selection vote per clause,
using the union over the six-point $\kappa$ grid. Both inculpatory
$S_\ell^+$ and exculpatory $S_\ell^-$ evidence must pass the stability
threshold; mean coefficient sign separates them. The report quarter
uses the fixed pool, signatures, and normalized score of~\eqref{eq:score}.
If the positive-weight denominator is zero, the implementation assigns
$s_\ell(x)=0$.}
\label{fig:cpss}
\end{figure*}

\subsection{Dataset Description}
We use CICIoMT2024~\cite{dadkhah2024}, WiFi/MQTT branch, with family-level
labels. The dataset was collected on a testbed of 40 healthcare devices
(25 real, 15 simulated) communicating over Wi-Fi, MQTT, and BLE. It
contains 18 attack types grouped into five attack families, and each flow is
described by 22 numeric CIC-extracted flow features. With benign traffic
there are six classes in all:
Benign, distributed denial-of-service (DDoS),
denial-of-service (DoS), attacks on the Message Queuing Telemetry Transport
protocol (MQTT), reconnaissance (Recon), and Spoofing.
\subsection{Booleanisation and Base Detector}
\label{sec:bool}
Each numeric flow feature is discretised into 10 quantile bins. We use
scikit-learn's KBinsDiscretizer with bin edges fit on the training
partition only. Features such as \texttt{Rate} contain $+\infty$ values
from division by near-zero flow duration upstream; we clip them to the
finite training maximum before binning, otherwise the quantile edges
collapse to a single degenerate bin (\S\ref{sec:validation}). The bins are
encoded as cumulative thermometer bits of the
form ``feature $\ge$ bin $k$''. TMU handles literal negation internally.
\subsection{Weighted TM}
\label{sec:wtm}
We train a weighted multiclass TM on the six
classes with $K = 400$ clauses per class, $T=320$, $s=8.0$, and 25 epochs.
Passing flows through the trained TM yields
the clause-activation matrix $Z \in \{0,1\}^{n \times CK}$ with $n$ flows
and $CK = 2{,}400$ clauses. Here
$Z[x,j]=1$ if and only if clause $j$ fires on input $x$.

\subsection{Clause-Pool Pruning}
\label{sec:prune}
Not every clause in the pool is worth carrying into signature extraction,
and which ones to drop can be posed exactly. Write $w_{c,j}$ for the
learned weight of clause $j$ of class $c$ and $z_{c,j}(x) \in \{0,1\}$ for
its activation, so the class sum is $f_c(x)=\sum_j w_{c,j} z_{c,j}(x)$ and
$\hat{y}(x)=\arg\max_c f_c(x)$. A keep-mask
$m \in \{0,1\}^{C \times K}$ with deletion sets $D_c=\{j: m_{c,j}=0\}$
gives $f_c(x;m)=\sum_j m_{c,j} w_{c,j} z_{c,j}(x)$, and we seek the
smallest pool that preserves accuracy:
\begin{equation}
\min_{m \in \{0,1\}^{C \times K}} \lVert m \rVert_0
\quad\text{s.t.}\quad
\mathrm{F_1}\!\left(\hat{y}(\cdot\,;m)\right) \ \ge\
\mathrm{F_1}\!\left(\hat{y}(\cdot\,;\mathbf{1})\right) - \varepsilon .
\label{eq:prune}
\end{equation}
This is a search over $2^{CK}$ masks ($2^{2400}$ here) whose constraint is
non-decomposable, macro-F1 sitting behind an $\arg\max$. We therefore
optimise the perturbation it induces,
$R(D_c) := \sum_{j \in D_c} |w_{c,j}|$, which controls~\eqref{eq:prune} through the
decision margin.

\begin{proposition}[Certificate]
\label{prop:cert}
Let $R = \max_c R(D_c)$ and let
$\gamma(x)=f_{\hat{y}(x)}(x)-\max_{c \neq \hat{y}(x)} f_c(x)$ be the
full-pool margin. If $\gamma(x) > 2R$ then $\hat{y}(x;m)=\hat{y}(x)$.
\end{proposition}
\begin{proof}
Since $z_{c,j}(x) \in \{0,1\}$,
$|f_c(x)-f_c(x;m)| = |\sum_{j \in D_c} w_{c,j} z_{c,j}(x)| \le R(D_c) \le R$
for every $c$. Writing $a=\hat{y}(x)$, for any $c \neq a$,
$f_a(x;m)-f_c(x;m) \ge \bigl(f_a(x)-R\bigr)-\bigl(f_c(x)+R\bigr)
\ge \gamma(x)-2R > 0$, so $a$ remains the maximiser.
\end{proof}

\begin{proposition}[Exact greedy]
\label{prop:greedy}
For a budget $|D_c| = k$, $\min_{D_c} R(D_c)$ is attained by the $k$
clauses of smallest $|w_{c,j}|$.
\end{proposition}
\begin{proof}
$R(D_c)$ is a sum of $k$ non-negative terms selected without replacement.
If $D_c$ contains $j$ while excluding $j'$ with $|w_{c,j'}| < |w_{c,j}|$,
swapping them strictly decreases $R$; no such swap exists precisely when
$D_c$ holds the $k$ smallest magnitudes.
\end{proof}

Magnitude ranking is therefore optimal for the surrogate rather than a
heuristic, and by Proposition~\ref{prop:cert} it carries a per-input
guarantee, at cost $O(CK\log K)$ against the $2^{CK}$ of~\eqref{eq:prune}. Sweeping
$k$ in steps of five and taking the largest $k$ that meets the constraint
gives the operating budget used in \S\ref{sec:pruned}; CAS is then
extracted from the surviving pool exactly as below, with nothing else
changed.

\subsection{CPSS Signature Extraction and Sign Filtering}
\label{sec:cpss}
CAS works entirely
in clause-activation space, where a flow is represented by the rules that
match it. We build one one-versus-rest problem
per family $\ell$. CPSS then repeats the following $B$ times: split the
fit half of the test set ($\mathrm{fit}$, \S\ref{sec:setup})
into two complementary halves, and on
each half fit an $\ell_1$-regularised logistic model from the columns of
$Z$ to the one-versus-rest label of family $\ell$ --- $2B$ half-fits per
family in all.

Two knobs control this stage: $\kappa$, the
inverse-regularisation strength of the logistic model inside each
half-fit, and $\pi_{\mathrm{thr}}$, the stability threshold applied after
all fits are done. A small $\kappa$ forces most coefficients to exactly
zero, so only a few clauses survive a single half-fit. We do not fix one
$\kappa$ but sweep it over a geometric
grid of 6 points from 0.001 to 0.1. A clause counts as selected by a
half-fit if any grid point selects it. This union-over-grid rule
follows~\cite{shah2013}.

Each half-fit casts one vote per clause. The stability score
$\hat{\pi}_j$ is the fraction of the $2B$ half-fits that select clause $j$.
The signature keeps the clauses that win enough votes:
$S_\ell = \{j : \hat{\pi}_j \ge \pi_{\mathrm{thr}}\}$. The knobs are not
independent: $\kappa$ sets which clauses can collect votes at all, and
$\pi_{\mathrm{thr}}$ then cuts the resulting vote distribution, so the
signature size at a given threshold depends on the grid. We restrict
$\kappa$ to the range that keeps each half-fit sparse on this
2{,}400-clause pool. A wider grid would inflate $q_\Lambda$ and make the
certificate vacuous.

One further step is needed. CPSS selects on $|$coefficient$|>0$, which is
correct for the error-control bound but wrong for attribution: a stably
negative coefficient is evidence \emph{against} $\ell$, not for it. We
therefore track the mean signed coefficient
$\overline{\mathrm{coef}}_j$ of each clause across fits and split the
signature into
$S_\ell^{+} = \{j \in S_\ell : \overline{\mathrm{coef}}_j > 0\}$
and $S_\ell^{-}$. The positive part is inculpatory evidence, the negative
part exculpatory. The certificate still applies to the full
sign-agnostic $S_\ell$; we state this scoping gap openly.
\S\ref{sec:validation} shows that this split is not cosmetic.

The attribution score is
\begin{equation}
s_\ell(x) = \frac{\sum_{j \in S_\ell^{+}} \hat{\pi}_j z_j(x) - \sum_{j \in S_\ell^{-}} \hat{\pi}_j z_j(x)}{\sum_{j \in S_\ell^{+}} \hat{\pi}_j},
\label{eq:score}
\end{equation}
and the prediction is $\hat{y} = \arg\max_\ell s_\ell(x)$. Matches to
inculpatory clauses add weight. Matches to exculpatory clauses subtract
weight. Normalisation by the inculpatory weight is required. Without it
scores scale with signature size instead of match quality.

\subsection{Error-Control Certificate}
\label{sec:cert}
With $q_\Lambda$ estimated as $\mathrm{mean}(\hat{\pi}) \cdot p$, the
Shah--Samworth bound gives
$E[V] \le q_\Lambda^2 / ((2\pi_{\mathrm{thr}} - 1) p)$ with $p = 2400$. We
report these bounds as loose wherever they are loose
(\S\ref{sec:signatures}): they certify that most of each signature is
likely genuine, not that it is near-noise-free.

\section{Experimental Setup}
\label{sec:setup}
Stratified caps are applied for tractability: at most 40{,}000 training
and 15{,}000 test rows per family, Benign (19{,}291) and Spoofing
(16{,}047) in full, giving 195{,}338 training and 80{,}868 test rows. The
distribution is heavily imbalanced (DDoS and DoS about 95.7\% of
flows~\cite{dadkhah2024}), so we report macro-F1 and per-class metrics.
The test set is permuted once and split three ways: CPSS half-fits use
$\mathrm{fit}$ ($40{,}434$ rows), every free choice --- the pruning budget
$k^{\star}$, the number of pairs $B$, and $\pi_{\mathrm{thr}}$ --- is made
on a disjoint selection quarter $\mathrm{sel}$ ($20{,}217$), and the report
quarter $\mathrm{rep}$ ($20{,}217$) is read exactly once, at the end. No
number we report was used to choose anything, which the closed-world
critiques of IDS evaluation ask
for~\cite{sommer2010,engelen2021,arp2022}. Multi-seed results use three
independently trained TMs (seeds 42, 7, 123) and fully seeded CPSS.

\section{Results}

\subsection{Validation of the Construction}
\label{sec:validation}
Three checks precede the headline. The infinite-value clipping of
\S\ref{sec:bool} is essential: without it macro-F1 collapses from 0.72
to 0.54. The sign split (\S\ref{sec:cpss}) is not cosmetic: the naive
sign-agnostic scorer --- a $\hat{\pi}$-weighted sum over the whole
selected support $S_\ell$ --- collapses to macro-F1 below $0.01$ on every
seed, while the signed scorer of~\eqref{eq:score} reaches $0.765$. And
$\pi_{\mathrm{thr}}{=}0.6$ is unstable as well as weak
($0.687 \pm 0.033$, one seed at $0.642$, against $0.743 \pm 0.016$ at
$\pi_{\mathrm{thr}}{=}0.7$), so we do not hand-pick a threshold: the
operating point is selected on $\mathrm{sel}$ and reported once on
$\mathrm{rep}$ (\S\ref{sec:setup}).

\subsection{Task A: Closed-Set Family Attribution}
\label{sec:taska}
At the operating point selected on $\mathrm{sel}$ ($B{=}15$,
$\pi_{\mathrm{thr}}{=}0.8$) CAS beats its base detector on every seed
(Table~\ref{tab:headline}); raising $B$ further adds nothing ($0.766$ at
$B{=}15$ against $0.751$ at $B{=}25$ on $\mathrm{sel}$). The gains land
where the base is weakest --- Benign recall rises from $0.26$ to $0.79$,
Spoofing precision from $0.43$ to $0.59$ --- while DDoS and DoS stay
mutually confused under both methods, a documented dataset
property~\cite{domenech2025}.

\begin{table}[!htb]
\centering
\caption{Task A: three-seed macro-F1 on the report quarter of
CICIoMT2024. The operating point $(B, \pi_{\mathrm{thr}})$ is selected on
the disjoint selection quarter; the pruned pool is \S\ref{sec:pruned}.}
\label{tab:headline}
\footnotesize
\begin{tabular}{@{}lcc@{}}
\toprule
Method & $(B, \pi_{\mathrm{thr}})$ & macro-F1 \\
\midrule
Weighted-TM argmax (base) & --- & $0.692 \pm 0.004$ \\
\textbf{CAS, full pool} & $(15, 0.8)$ & $\mathbf{0.765 \pm 0.029}$ \\
\textbf{CAS, pruned pool} & $(10, 0.9)$ & $\mathbf{0.787 \pm 0.023}$ \\
\bottomrule
\end{tabular}
\end{table}

\subsection{Sensitivity to $\pi_{\mathrm{thr}}$, $B$, and $\kappa$}
\label{sec:sweep}
Fig.~\ref{fig:sweep} sweeps both knobs on the report quarter after
pruning to $390$ clauses ($65$ per class). At $B{=}15$, varying
$\pi_{\mathrm{thr}}$ from $0.5$ to $0.9$ leaves macro-F1 within
$0.777$--$0.779$ and accuracy within $0.792$--$0.794$. At
$\pi_{\mathrm{thr}}{=}0.8$, varying $B$ from $5$ to $25$ leaves
macro-F1 within $0.776$--$0.778$ and accuracy within $0.790$--$0.793$.
Both metrics remain above the pruned weighted-TM baselines
($0.688$ macro-F1 and $0.720$ accuracy). These sensitivity slices
retain the original fixed settings for comparison; the pruned-pool
operating point selected on the selection quarter remains
$(B,\pi_{\mathrm{thr}})=(10,0.9)$.
The separate $\kappa$ sweep uses the full $2{,}400$-clause pool
(Table~\ref{tab:grid}): across
$[10^{-3}, 4\times10^{-2}]$ the mean stays within $0.724$--$0.763$,
degrading only at $\kappa{=}0.1$. This justifies the union-over-grid rule:
at the selected threshold the union ($0.765$) matches or exceeds every
single-$\kappa$ row (best $0.763$).

\begin{figure}[!htb]
\centering
\includegraphics[width=0.88\columnwidth]{fig_sweep_pruned.pdf}
\caption{Sensitivity after pruning to $390$ clauses ($65$ per class),
on the report quarter, three-seed mean$\pm$sd.
(a)~CAS macro-F1 and accuracy versus $\pi_{\mathrm{thr}}$ at $B{=}15$;
dotted lines are the pruned weighted-TM baselines. (b)~versus the number of
complementary pairs $B$ at $\pi_{\mathrm{thr}}{=}0.8$.}
\label{fig:sweep}
\end{figure}

\begin{table}[!htb]
\centering
\caption{$\kappa \times \pi_{\mathrm{thr}}$: three-seed mean CAS macro-F1
on the report quarter, $B{=}15$.
The bottom row is the union-over-grid rule actually used, which matches or
exceeds every single $\kappa$ at the selected $\pi_{\mathrm{thr}}{=}0.8$.}
\label{tab:grid}
\footnotesize
\begin{tabular}{@{}lccccc@{}}
\toprule
$\kappa \backslash\ \pi_{\mathrm{thr}}$ & 0.5 & 0.6 & 0.7 & 0.8 & 0.9 \\
\midrule
$1.0\!\times\!10^{-3}$ & 0.728 & 0.725 & 0.724 & 0.724 & 0.727 \\
$2.5\!\times\!10^{-3}$ & 0.748 & 0.745 & 0.739 & 0.741 & 0.746 \\
$6.3\!\times\!10^{-3}$ & 0.751 & 0.746 & 0.737 & 0.740 & 0.763 \\
$1.6\!\times\!10^{-2}$ & 0.748 & 0.742 & 0.752 & 0.752 & 0.745 \\
$4.0\!\times\!10^{-2}$ & 0.743 & 0.752 & 0.747 & 0.742 & 0.755 \\
$1.0\!\times\!10^{-1}$ & 0.673 & 0.693 & 0.718 & 0.722 & 0.694 \\
\midrule
union grid (method) & 0.708 & 0.687 & 0.743 & $\mathbf{0.765}$ & 0.762 \\
\bottomrule
\end{tabular}
\end{table}

\subsection{Binary Reduction: The Attack Signature}
\label{sec:binary}
Collapsing the five attack families into one \textsc{attack} class and
training a dedicated two-class TM yields a binary CAS on the same
booleanisation. The base TM
over-predicts attack, catching almost everything (recall 0.987) but crying
wolf often (precision 0.844). The attack signature ($B{=}25$,
$\pi_{\mathrm{thr}}{=}0.9$, selected on $\mathrm{sel}$) rebalances the
operating point (recall 0.910, precision 0.946) and ranks flows far
better, lifting AUROC from 0.817 to 0.938 (Table~\ref{tab:binary}).

\begin{table}[!htb]
\centering
\caption{Binary attack detection on the report quarter, three seeds
(mean).}
\label{tab:binary}
\footnotesize
\begin{tabular}{@{}lccc@{}}
\toprule
Method & attack F1 & R / P & AUROC \\
\midrule
Binary TM argmax & $0.910 \pm 0.000$ & 0.987 / 0.844 & $0.817 \pm 0.045$ \\
\textbf{CAS signature} & $\mathbf{0.928 \pm 0.008}$ & 0.910 / 0.946 & $\mathbf{0.938 \pm 0.006}$ \\
\bottomrule
\end{tabular}
\end{table}

\subsection{Signature Content and Certificates}
\label{sec:signatures}
At $\pi_{\mathrm{thr}}{=}0.6$, $|S|$ ranges 115--311, $|S^{+}|$ 29--58,
and $E[V]$ ${\le}72$--${\le}250$. At the selected point
$\pi_{\mathrm{thr}}{=}0.8$
(Table~\ref{tab:reduction}, left) each family shrinks to 28--53
sign-positive clauses while F1 \emph{improves}. The decoded DDoS signature
is the classic flooding profile (\texttt{Rate} $\ge$ 31.5k packets/s,
short headers, ICMP tests), while exculpatory lists are long low-rate
conjunctions resembling benign chatter.


\subsection{Clause-Pool Pruning and CAS on the Pruned Pool}
\label{sec:pruned}
Solving the pruning program of \S\ref{sec:prune},
Fig.~\ref{fig:pruning} traces macro-F1 on the selection quarter as $k$
grows in steps of five: flat
until $335$ of the $400$ clauses per class ($84\%$) are gone, so
$\varepsilon = 0.01$ gives $k^{\star} = 335$ without touching any
evaluation data. The surrogate certifies but
does not solve~\eqref{eq:prune}: a global knapsack on $\sum_c R(D_c)$, or
the minimax rule equalising $\max_c R(D_c)$, changes macro-F1 by less than
the cross-seed spread ($0.689$, $0.678$, $0.683$ at $390$ clauses).

\begin{figure}[!htb]
\centering
\includegraphics[width=0.48\columnwidth]{clause_pruning.png}
\caption{Clause-weight pruning on the trained weighted TMs, three-seed
mean$\pm$sd on the report quarter. Clauses are removed per class in order
of increasing $|$weight$|$, five at a time. Macro-F1 and accuracy are flat
until $84\%$ of the pool has been deleted.}
\label{fig:pruning}
\end{figure}

If $84\%$ of the pool is redundant for the base detector, is it also
diluting the stability vote? Keeping only the top $65$ clauses per class
--- the knee of Fig.~\ref{fig:pruning} --- and re-running the identical
pipeline on the $390$-clause pool answers yes (Fig.~\ref{fig:prunedcas}). Pruning costs the base detector nothing
($0.688 \pm 0.008$ vs.\ $0.692 \pm 0.004$), but CAS reaches
$0.787 \pm 0.023$ at its selected point ($B{=}10$,
$\pi_{\mathrm{thr}}{=}0.9$), above the $0.765 \pm 0.029$ from the full pool. More
useful is the effect on the threshold: the full-pool sweep swings from
$0.687$ to $0.765$ with $\pi_{\mathrm{thr}}$, whereas the pruned pool is
flat at $0.775$--$0.787$ throughout, so the instability of
\S\ref{sec:validation} is a property of the diluted pool, not of CPSS.
Signatures compact correspondingly (Table~\ref{tab:reduction}, right):
cores of 22--42 clauses from a six-times-smaller pool, $E[V]$ tightening to
${\le}27$--${\le}50$.


\begin{table}[!htb]
\centering
\caption{Signature compaction at $\pi_{\mathrm{thr}}{=}0.8$ (seed 42) on
the full $2{,}400$-clause pool and on the weight-pruned $390$-clause pool.}
\label{tab:reduction}
\footnotesize
\begin{tabular}{@{}lcccc@{}}
\toprule
& \multicolumn{2}{c}{full pool ($2{,}400$)} & \multicolumn{2}{c}{pruned ($390$)} \\
\cmidrule(lr){2-3}\cmidrule(lr){4-5}
Family & $|S|$ & $|S^{+}|$ & $|S|$ & $|S^{+}|$ \\
\midrule
Benign   & 184 & 41 & 63 & 30 \\
DDoS     &  63 & 29 & 66 & 36 \\
DoS      & 247 & 33 & 83 & 42 \\
MQTT     & 145 & 53 & 64 & 26 \\
Recon    &  66 & 28 & 58 & 22 \\
Spoofing & 194 & 50 & 76 & 28 \\
\bottomrule
\end{tabular}
\end{table}

\begin{figure}[!htb]
\centering
\includegraphics[width=0.48\columnwidth]{pruned_vs_full.png}
\caption{CAS macro-F1 versus $\pi_{\mathrm{thr}}$ on the full and
weight-pruned clause pools, three-seed mean$\pm$sd on the report quarter
at $B{=}15$. Pruning lifts the
curve and flattens it: on the pruned pool the stability threshold no
longer changes the result.}
\label{fig:prunedcas}
\end{figure}

\subsection{Open-Set Novelty}
\label{sec:taskb}
Leave-one-family-out with split-conformal thresholds ($\epsilon{=}0.05$,
two folds) is negative: the false-flag guarantee holds (FPR 0.035 Spoofing,
0.027 MQTT, both under target) but power is weak (Spoofing AUROC 0.275,
below chance; MQTT 0.622) --- ARP-spoofing closely resembles legitimate ARP
discovery in this feature space, a measured blind spot.

\subsection{Mixture Identifiability}
\label{sec:disjunct}
With the six
inculpatory signatures as columns of a binary clause-by-family table, the
largest $t$ for which no column is covered by the union of any $t$ others
is $\hat{t}{=}5$ --- the maximum for six families --- on every seed and
threshold, so any noiseless mixture of up to five families decodes exactly. Under $5\%$ noise (500 trials) strict COMP collapses to
$0.2\%$--$1.8\%$, while a partial-match decoder restores $99.8\%$--$100\%$
at threshold 0.8 and $68\%$--$76\%$ at $15\%$ noise.

\subsection{Cross-Seed Signature Stability}
\label{sec:stability}
The pairwise Kuncheva index on clause-index overlap across seeds
is near zero for every family ($-0.015$ to $0.111$): CAS stabilises
selection \emph{within} a trained TM --- what attribution relies on ---
but signature identity does not survive retraining.

\section{Discussion and Conclusion}
\label{sec:discussion}
CAS turns a trained TM's clause pool into compact signed evidence lists
per family, with a bound on spurious inclusions and exact mixture
identifiability. On CICIoMT2024 it beats its base detector (macro-F1 0.765
vs.\ 0.692; binary AUROC 0.938), and discarding $84\%$ of the pool by
weight first raises this to 0.787 and removes the threshold sensitivity.
All operating points are chosen on a disjoint selection quarter
(\S\ref{sec:setup}), so no reported number was used to choose anything.
Limitations: no post-hoc XAI baseline, loose $E[V]$ bounds,
near-chance cross-seed identity, single-dataset scope.

\begin{thebibliography}{15}
\bibitem{granmo2018} O.-C. Granmo, ``The Tsetlin Machine --- a
game-theoretic bandit-driven approach to optimal pattern recognition with
propositional logic,'' \emph{arXiv:1804.01508}, 2018.
\bibitem{phoulady2020} O.-C. Granmo and S. Phoulady, ``The Weighted Tsetlin
Machine,'' \emph{arXiv:2002.01245}, 2020.
\bibitem{jaiswal2026iomt} R. Jaiswal, P.-A. Andersen, L. R. Cenkeramaddi,
L. Jiao, and O.-C. Granmo, ``A Tsetlin Machine-driven Intrusion Detection
System for Next-Generation IoMT Security,'' \emph{IEEE SVCC} /
\emph{arXiv:2604.03205}, 2026.
\bibitem{shah2013} R. Shah and R. Samworth, ``Variable selection with error
control: another look at stability selection,'' \emph{JRSS-B}, vol. 75,
no. 1, pp. 55--80, 2013.
\bibitem{tmu} O.-C. Granmo et al., ``TMU: Tsetlin Machine Unified,''
github.com/cair/tmu.
\bibitem{abeyrathna2020} K. D. Abeyrathna et al., ``Intrusion Detection
with Interpretable Rules Generated using the Tsetlin Machine,'' \emph{IEEE
SSCI}, 2020.
\bibitem{jaiswal2026medsec} R. Jaiswal et al., ``On-Device Interpretable
Tsetlin Machine-Based Intrusion Detection for Secure IoMT,''
\emph{arXiv:2605.16707}, 2026.
\bibitem{meinshausen2010} N. Meinshausen and P. B\"uhlmann, ``Stability
selection,'' \emph{JRSS-B}, vol. 72, no. 4, pp. 417--473, 2010.
\bibitem{sommer2010} R. Sommer and V. Paxson, ``Outside the Closed World: On
Using Machine Learning for Network Intrusion Detection,'' \emph{IEEE S\&P},
2010.
\bibitem{engelen2021} G. Engelen, V. Rimmer, and W. Joosen,
``Troubleshooting an Intrusion Detection Dataset: the CICIDS2017 Case
Study,'' \emph{IEEE SPW}, 2021.
\bibitem{arp2022} D. Arp et al., ``Dos and Don'ts of Machine Learning in
Computer Security,'' \emph{USENIX Security}, 2022.
\bibitem{dadkhah2024} S. Dadkhah et al., ``CICIoMT2024: A benchmark dataset
for multi-protocol security assessment in IoMT,'' \emph{Internet of Things},
vol. 28, art. 101351, 2024.
\bibitem{domenech2025} J. Dom\'enech Fons, ``Evaluating and enhancing
intrusion detection systems in IoMT: The importance of domain-specific
datasets,'' \emph{Internet of Things}, 2025.
\bibitem{xaiids2025} ``Explainable AI for Interpreting ML-based IDS using
LIME and SHAP,'' \emph{Applied Sciences}, vol. 15, no. 13, 2025.
\bibitem{du1993} D.-Z. Du and F. K. Hwang, \emph{Combinatorial Group Testing
and Its Applications}, World Scientific, 1993.
\end{thebibliography}

\end{document}
